\chapter{Topology}

\section{Topological Spaces}

\begin{definition}
    A topological space is a tuple \((X,\mca{T})\),
    where \(X\) is a set, \(\mca{T} \subseteq \mca{P}(X)\) a collection of subsets,
    such that:
    \begin{enumerate}[label={(\roman*)}, itemsep=0mm]
        \item \(\emptyset,X \in \mca{T}\);
        \item closed under finite intersections
            \(U,V \in \mca{T} \implies U \cap V \in \mca{T}\); and
        \item closed under arbitrary unions
            \({\{U_\alpha\}} \subset \mca{T} \implies \bigcup_\alpha U_\alpha \in \mca{T}\).
    \end{enumerate}
    \(\mca{T}\) is called a topology on the set \(X\);
    elements in \(\mca{T}\) are called open,
    while the complement of these elements are called closed.
\end{definition}
\begin{definition}
    Suppose \(X\) is a topological space.
    Given any \(x \in X\), a neighbourhood of \(x\) is a set \(V \subseteq X\)
    such that there exists an open set \(U\) with \(x \in U \subseteq V\).
\end{definition}
\begin{proposition}
    Suppose \(X\) is a topological space.
    \(U \subseteq X\) is open if and only if \(U\) is a neighbourhood of all \(x \in U\).
\end{proposition}
\begin{proof}
    Suppose \(U\) is open. Then it is trivially a neighbourhood for all \(x \in U\).

    Suppose \(U\) is a neighbourhood for all \(x \in U\).
    Then for each \(x\), there exists an open set \(U_x\) such that \(x \in U_x \subseteq U\).
    But then \(\bigcup_{x \in U} U_x \supseteq U\) because it contains all points of \(U\);
    and the other inclusion also holds because each of the \(U_x \subseteq U\).
\end{proof}

\subsection*{Separation Axioms}

\begin{remark}
    Despite the name `separation axioms',
    these are a set of useful topological definitions, rather than axioms.
\end{remark}

\begin{definition}
    Suppose \(X\) is a topological space, and \(x,y \in X\) are distinct points (as in \(x \neq y\)).
    \(x,y\) are topologically distinguishable if there exists an open set \(U \subseteq X\)
    such that either \(x \in U\) and \(y \notin U\) or \(y \in U\) and \(x \notin U\).
    A topological space is Kolmogorov or \(\text{T}_0\) if every pair of points is topologically distinguishable.
\end{definition}

\begin{definition}
    Suppose \(X\) is a topological space, and \(x,y \in X\) are distinct points.
    \(x,y\) are separated if there exists neighbourhoods \(U \ni x\) and \(V \ni y\)
    such that \(x \notin V\) and \(y \notin U\).
    A topological space is Fr\'echet or \(\text{T}_1\) if every pair of points is separated.
\end{definition}
\begin{lemma}
    \(\text{T}_1\) spaces are \(\text{T}_0\).
\end{lemma}
\begin{proof}
    Suffices to prove that separated points are topologically distinguishable.
    If \(x,y\) are separated, then there exists some neighbourhood \(U \ni x\) with \(y \notin U\),
    which gives us topological distinguishability.
\end{proof}
\begin{proposition}
    A topological space is \(\text{T}_1\) if and only if every singleton set \(\{x\} \subseteq X\) is closed.
\end{proposition}
\begin{proof}
    Suppose \(X\) is \(\text{T}_1\), and fix an arbitrary \(x \in X\).
    For every \(y \neq x\), there exists a neighbourhood \(V_y \ni y\) such that \(x \notin V_y\).
    By definition of a neighbourhood, for each of these \(y\),
    there exists an open \(U_y\) such that \(y \in U_y \subseteq V_y\).
    But \(\bigcup_{y \neq x} U_y\) is open, and contains every point but \(x\),
    which means it is \({\{x\}}^c\).
    Then \(\{x\}\) is closed.

    Now, suppose every singleton is closed.
    Picking any two distinct \(x,y \in X\), we have \(y \in {\{x\}}^c\) open and \(x \in {\{y\}}^c\) open.
\end{proof}

\begin{definition}
    A topological space \(X\) is called Hausdorff or \(\text{T}_2\)
    if any two points \(x,y \in X\) have neighbourhoods \(U \ni x\) and \(V \ni y\)
    such that they are disjoint (\(U \cap V = \emptyset\)).
\end{definition}
\begin{lemma}
    \(\text{T}_2\) spaces are \(\text{T}_1\).
\end{lemma}
\begin{proof}
    Since the neighbourhoods are disjoint, in particular, they do not contain the other point.
\end{proof}

\begin{definition}
    A topological space \(X\) is called regular if for all \(x \in X\) and all closed sets \(C \subseteq {\{x\}}^c\),
    there exists disjoint open sets \(U \ni x\) and \(V \supseteq C\) such that \(U \cap V = \emptyset\).
    A regular Hausdorff space is called \(\text{T}_3\).
\end{definition}

\begin{definition}
    A topological space \(X\) is called normal if for all disjoint closed subsets \(C,D \subseteq X\),
    there exists open sets \(U \supseteq C\) and \(V \supseteq D\) such that \(U \cap V = \emptyset\).
    A normal Hausdorff space is called \(\text{T}_4\).
\end{definition}

\subsection*{Continuity}

\begin{definition}
    Suppose \(\func{f}{X}{Y}\) is a function between topological spaces.
    \(f\) is continuous if the preimage of any open set is open.
\end{definition}
\begin{definition}
    A continuous function \(\func{f}{X}{Y}\) with a continuous inverse \(\func{f^{-1}}{Y}{X}\) is a homeomorphism.
\end{definition}

\subsection*{Bases}

\begin{definition}
    Suppose \(X\) is a topological space, and \(x \in X\) is a point.
    A collection \(\mca{B}_x\) of open neighbourhoods of \(x\) is a local basis at \(x\)
    if for all open \(U \ni x\), there exists some \(B \in \mca{B}_x\) such that \(x \in B \subseteq U\).
    If every point admits a countable local basis,
    then we say \(X\) is first countable.
\end{definition}
% \begin{definition}
%     A basis \(\mca{B}\) is a collection of subsets of \(X\) such that:
%     \begin{enumerate}[label={(\roman*)}, itemsep=0mm]
%         \item for every point \(x \in X\), there exists some basis element \(B \in \mca{B}\) such that \(x \in B\); and
%         \item if \(x \in B_1 \cap B_2\) where \(B_1,B_2 \in \mca{B}\),
%             then there exists some basis element \(B_3 \in \mca{B}\) such that \(x \in B_3 \subseteq B_1 \cap B_2\).
%     \end{enumerate}
%     The topology \(\mca{T}\) generated by the basis \(\mca{B}\)
%     is the collection of all sets \(U\) (which we define to be open)
%     such that for each \(x \in U\), there exists a basis element \(B \in \mca{B}\) such that \(x \in B \subseteq U\).
% \end{definition}
\begin{definition}
    Suppose \(X\) is a topological space.
    A collection \(\mca{B}\) of open sets is a basis if it contains a local basis for every point.
    If \(\mca{B}\) has a countable basis, then we say \(X\) is second countable.
\end{definition}
% \begin{lemma}
%     Suppose \(\mca{B}\) is a basis, and \(\mca{T}\) is the topology generated by \(\mca{B}\).
%     Then:
%     \begin{enumerate}[label={(\alph*)}, itemsep=0mm]
%         \item \(\mca{T}\) is the set of all arbitrary unions of elements of \(\mca{B}\); and
%         \item \(\mca{T}\) indeed forms a topology.
%     \end{enumerate}
% \end{lemma}
% \begin{proof}
%     First off, clearly any arbitrary unions of \(\mca{B}\)-sets contain \(\mca{B}\)-sets,
%     so they are contained within \(\mca{T}\).
% \end{proof}
% \begin{lemma}
%     Suppose \(\mca{T}\) is a topology on \(X\)
%     and \(\mca{B} \subseteq \mca{T}\) is a collection of open sets
%     such that given any open set \(U\), for each \(x \in U\),
%     there exists some \(B \in \mca{B}\) with \(x \in B \subseteq U\).
%     Then \(\mca{B}\) is a basis for the topology \(\mca{T}\).
% \end{lemma}

\subsection*{Induced Topologies}

\begin{definition}
    Suppose \({\{X_i\}}_{i \in I}\) is a family of topological spaces, and \(Y\) is a set.
    If we have functions \(\func{f_i}{Y}{X_i}\) for each \(i\),
    the induced topology on \(Y\) is the topology with fewest open sets
    such that \(f_i\) are all continuous.
\end{definition}

\begin{definition}
    Suppose \(X\) is a topological space.
    If \(Y \subseteq X\) is a subset, then the set of all \(U \cap Y\) where \(U \subseteq X\) is open
    forms the subspace topology for \(Y\).
\end{definition}
\begin{proposition}[Universal Property of Subspace Topology]
    Suppose \(X\) is a topological space, and \(Y \subseteq X\) is a subset.
    % and \(\) is the inclusion map.
    Then the subspace topology on \(Y\) is the topology induced by the inclusion map \(\func{\iota}{Y}{X}\).
\end{proposition}
\begin{proof}
    It is obvious that if we remove any open set \(U \cap Y \subseteq X\),
    then the preimage \(\iota^{-1}(U)\) is no longer open.
\end{proof}

\begin{definition}
    Suppose \({\{X_i\}}_{i \in I}\) are topological spaces.
    Then the product topology of \(\prod_{i \in I} X_i\)
    % is the topology generated by (the basis containing) the sets \(U \times V\)
    % for all open \(U \subseteq X\) and \(V \subseteq Y\).
    is the topology induced by the projection maps \(\func{f_i}{\prod_{i \in I} X_i}{X_i}\).
\end{definition}

\begin{theorem}
    Suppose \({\{X_i\}}_{i \in I}\) are topological spaces.
    If \(A_i \subseteq X_i\) are subspaces,
    then the product topology on \(\prod A_i\) is identical to the subspace topology of \(\prod A_i \subseteq \prod X_i\).
\end{theorem}
\begin{proof}
    It is obvious that we have the commutative diagram
    \begin{equation*}
        \begin{tikzcd}
            \prod A_i \arrow[hookrightarrow]{r} \arrow[twoheadrightarrow]{d} & \prod X_i \arrow[twoheadrightarrow]{d} \\
            A_i \arrow[hookrightarrow]{r} & X_i
        \end{tikzcd}
    \end{equation*}
    so both topologies are the topology induced by the maps \(\prod A_i \to X_i\).
\end{proof}

\subsection*{Coinduced Topologies}

\begin{definition}
    Suppose \({\{X_i\}}_{i \in I}\) is a family of topological spaces, and \(Y\) is a set.
    If we have functions \(\func{f_i}{X_i}{Y}\) for each \(i\),
    the coinduced topology on \(Y\) is the topology with the most open sets
    such that \(f_i\) are all continuous.
\end{definition}
\begin{definition}
    Suppose \(X\) is a topological space and \(\func{f}{X}{Y}\) is a surjective function.
    The quotient topology on \(Y\) is the coinduced topology by \(f\).
\end{definition}


\section{Compactness}

\begin{definition}
    A topological space \(X\) is considered quasi-compact
    if every open cover of \(X\) has a finite subcover.
    A topological space \(X\) is considered compact
    if it is quasi-compact and Hausdorff.
\end{definition}
